\section{模型假设}

题目已经给出储能设备的主要参数，但对售电机制、日前信息边界和跨日终端状态等细节未全部规定。以下假设只用于补足这些必要条件，并尽量不额外引入题面之外的市场机制。
\begin{enumerate}[label=\arabic*., itemsep=2pt, topsep=3pt]
    \item \textbf{统一时间尺度。}以 $10$ min 为一个调度时段，一天共有 $T=144$ 个时段，$\Delta t=1/6$ h；附件中的 $0{:}00{+}1$ 视为当天 $24{:}00$。附件各行与结果模板按原顺序对应，不再人为平移时间标签。
    \item \textbf{不考虑反向售电。}题目只规定从外网购电，没有给出向外网售电的价格和交易规则，外网购电量据此保持非负；光伏在满足负载和储能充电后仍无法利用的部分记为弃光。
    \item \textbf{储能参数全年不变。}主模型取充电效率和放电效率均为 $0.9$，忽略自放电、容量衰减、循环寿命成本和启停损耗。由于“充放电效率为 $90\%$”存在解释空间，另一种往返效率定义在灵敏度分析中另行检验。
    \item \textbf{功率上限按时段折算。}最大充、放电功率 $5000$ kW 在一个 $10$ min 时段内对应最大充、放电量 $5000\Delta t=833.33$ kWh，所有充放电约束均按该电量上限执行。
    \item \textbf{全年储能状态连续。}2025年1月1日 $0{:}00$ 储电量为 $6000$ kWh，随后每日实际日末储电量直接作为下一日初值。官方结果从2月1日起填写，1月既用于历史预测冷启动，也用于把储能状态自然递推到正式评价区间。
    \item \textbf{问题二只使用决策时刻之前的历史数据。}在 $0{:}00$ 制定全天计划时，不直接使用当天尚未实现的真实负载和真实光伏，而由历史数据形成因果预测；“日前负荷完全已知”仅作为信息更充分的对照，用来衡量信息价值。
    \item \textbf{冷启动使用附件1先验。}1月1日前不存在历史观测时，以附件1的典型日负载和光伏曲线作为一次性初始先验；一旦积累真实历史数据，即转入滚动历史预测，不再长期依赖附件1源荷曲线。
    \item \textbf{允许储能实时缓冲计划偏差。}问题二的全天计划购电量在日内锁定，但当期真实负载、光伏和储电量实现后，可以调整储能充放电。该调整只使用当前及过去信息，不读取未来真实数据；储能仍无法覆盖的供电缺口才计为紧急购电。
    \item \textbf{问题二、三计划层采用日循环终端状态。}除问题一的题面硬约束外，问题二、三主模型令计划终端储电量回到当日实际初始储电量，以避免有限时域优化在午夜通过透支电池获得虚假的短期收益；取消该约束的情形在敏感性分析中单独比较。
    \item \textbf{问题四主结果直接使用附件4电价。}这是题面明确给出的重算数据；仅用历史价格预测未来价格的方案作为现实市场信息受限时的扩展分析，不替代正式结果。
\end{enumerate}

\section{符号说明}

主要符号列于表~\ref{tab:symbols}。为避免功率与电量混淆，负载和光伏仍以 kW 表示，进入每个 $10$ min 时段的能量平衡前乘以 $\Delta t$；购电、充放电、弃光和紧急购电均直接以 kWh 表示。$S_{d,t}$ 表示储能内部实际储电量，而非百分比形式的 SOC。

\begin{table}[htbp]
    \caption{主要符号说明}
    \label{tab:symbols}
    \centering
    \zihao{-5}
    \begin{tabularx}{0.94\textwidth}{>{\centering\arraybackslash}p{2.5cm} >{\centering\arraybackslash}p{1.8cm} X}
        \toprule
        \multicolumn{1}{c}{符号} & \multicolumn{1}{c}{单位} & \multicolumn{1}{c}{含义} \\
        \midrule
        $d,\ t$ & --- & 日期下标与日内 $10$ min 时段下标（$t=1,\dots,144$）\\
        $\Delta t$ & h & 单个调度时段长度，$\Delta t=1/6$ \\
        $L_{d,t}$ & kW & 小区负载功率 \\
        $P_{d,t}$ & kW & 光伏实际发电功率 \\
        $\hat P_{d,t}$ & kW & 决策时刻可用的光伏预测功率 \\
        $\pi_{d,t}$ & 元/kWh & 外网正常购电单价 \\
        $G^{p}_{d,t}$ & kWh & $0{:}00$ 制定的计划购电量 \\
        $G^{a}_{d,t}$ & kWh & 执行前最后一次有效调整后的购电量 \\
        $C_{d,t}$ & kWh & 流入储能的充电量 \\
        $H_{d,t}$ & kWh & 储能输出的放电量 \\
        $S_{d,t}$ & kWh & 第 $t$ 时段结束时储能内部储电量 \\
        $W_{d,t}$ & kWh & 未被利用的弃光电量 \\
        $E_{d,t}$ & kWh & 供电不足时的紧急购电量 \\
        $u_{d,t},\ v_{d,t}$ & kWh & 调整购电量相对计划的上调量 $(G^a-G^p)^+$ 与下调量 $(G^p-G^a)^+$ \\
        $z_{d,t}$ & --- & 充放电互斥变量；$z=1$ 允许充电，$z=0$ 允许放电 \\
        $\eta_c,\ \eta_d$ & --- & 充电效率与放电效率，主模型均取 $0.9$ \\
        $S_{\min},\ S_{\max}$ & kWh & 储电量上下限，分别为 $1200$ 与 $10800$ \\
        $\bar E$ & kWh & 单时段充放电量上限，$\bar E=5000\Delta t=833.33$ \\
        $\mathcal{I}_{d,\tau}$ & --- & 第 $d$ 天决策节点 $\tau$ 已经获得的信息集合 \\
        \bottomrule
    \end{tabularx}
\end{table}

后文最容易混淆的是 $G^p$、$G^a$ 与 $E$。$G^p$ 是 $0{:}00$ 对未来各时段作出的基准承诺，$G^a$ 是问题三中在日内新预报到达后、交易时段执行前最后一次生效的调整量；二者的差额决定调整结算。即使已经进行过调整，真实运行仍可能再次偏离预测，因此还保留紧急购电 $E$ 作为最终供电兜底。三个量分别对应“日初计划—日内修订—实际缺口”，其时间顺序贯穿问题二、三的全部建模。
